\cotefile{78}
\scanpage{61}

Pour prouver \unsure{$\Leftarrow$)|$\Leftarrow$|a)}, notons que $\sigma_{i}R^+=R^-$ $\Longleftrightarrow \sigma_i R^-=R^+\unsure{,|{}}$ $(i\in\{0,\unsure{2|1}\})$, donc si $R^+$ est une orientation, on aura $uR^+=R^+$ donc $G_1^+(R^+)\subset R^+$, et $gR^+=R^-$ si $g\in G_1\setminus G_1^0=G_1^-$, \deleted{Il suffit de voir que} et inversement. Or dans le cas d'un polygone, \deleted{\unsure{$\simeq G$}} \unsure{donc} $\simeq\Pi_{n}$, on voit immédiatement que $G_1^+$ opérant sur $R\simeq\mathbb D_n$ a \inserted{\unsure{exactemt|exactement}} deux orbites, \deleted{qui} \inserted{lesquelles sont} échangées par $\sigma$, savoir $D_n^+$ et $D_n^-$, -- ce sont donc \unsure{celà|là|ici} les orientations.

\deleted{\illegible{1}} \textcircled{7}. \uline{Prop} Soit $\Pi$ \inserted{$=(S,A,R)$} polygone combinatoire, \deleted{\unsure{$G$}} \inserted{désignons} $G=\operatorname{Aut}\Pi$, \deleted{\unsure{On|donc}} $\operatorname{Or}(\Pi)$ a deux éléments et \deleted{\unsure{pour|par|$G$}} transport de structure \dots $G\longrightarrow\unsure{\mathfrak S|\mathscr S}_{\operatorname{Or}(\Pi)}\simeq\pm1$, \deleted{\unsure{Rappelons|que l'on}} Soit $G^+$ le noyau, a) On \deleted{\unsure{obt}} \inserted{obtient ainsi} suite exacte

\[
1\longrightarrow G^+\longrightarrow G\mathbin{-}\{\pm1\}\mathbin{-}1
\]

et tt élément de $G^-=G\setminus G^+$ est d'ordre $2$, donc splitte cette suite exacte. \deleted{b) $R$ est torseur} \deleted{\unsure{a)}} $G^+$ est cyclique d'ordre $n$, \deleted{si $R^+\in\operatorname{Or}(\Pi)$, et $S,A$ \unsure{sont des}} \inserted{et l'él. $-1$ de $\pm1$} y opère par symétrie.

b) $R$ est un torseur sous $G$, $S,A$ sont torseurs sous $G^+$.

% archive sheet number: 60
